\section{Introduction}
\label{sec:introduction}

An assistant that remembers creates a durable information system. A sentence accepted today may be recalled months later, combined with other facts, shown to a model, or used while proposing a change to the outside world. The engineering question is therefore not only ``did retrieval find a relevant sentence?'' It is also: who said it, was that source trusted, what did it replace, can it be deleted, why did it surface, and does it grant authority to act?

Most of those questions must be answerable \emph{after the fact}, from evidence, by someone who does not trust the assistant's own account of events. We call that capability \emph{recall forensics} (\cref{sec:forensics}) and treat it as a first-class design goal rather than a logging afterthought: each measured recall records a bounded candidate ledger and returned graph paths, the records they cite chain back to source commitments, and selected fields are digest-bound to a per-subject hash chain. A standard-library auditor can therefore reconstruct the measured decisions and detect the mutation classes evaluated in this paper, subject to explicit deletion and candidate-generation limits.

Many memory implementations begin with a document store and semantic similarity. That is useful for search, but search alone does not define durable state. Similarity does not decide whether a web page may overwrite a user's preference; a vector does not itself express correction, quarantine, deletion, approval, or an uncertain external result. \system{} starts from these state-transition questions and treats retrieval as one subsystem inside a larger control plane.

\plainbox{Think of the system as four ordinary things used together: a filing cabinet for the original statements, a card catalogue for finding them, a map connecting related people and facts, and a receipt book recording every important change. The map and catalogue can be rebuilt; the original statements and receipts are the durable evidence.}

This paper documents the core architecture implemented in the open-source \system{} release v0.3.0, at commit \code{69ebedea93ea95965867af5483dd9f63425b8e75}~\cite{aetnamem2026}. The recall-forensics bindings and evaluator are a later, content-addressed development snapshot based on commit \code{12925d3}; the evidence manifest, rather than that base commit alone, identifies its exact source bytes. The paper makes six scoped contributions.

\begin{enumerate}
  \item It defines \textbf{recall forensics} as five scoped questions --- attribution, provenance commitment, lifecycle consistency, decision replay, and integrity --- and evaluates them with a standard-library auditor: 14/14 rankings and 68/68 returned-result pairs reconstructed, 11/11 eligible turns re-executed exactly, 75/75 selected single-row mutations detected, and three privacy-affected turns explicitly excluded from engine re-execution.
  \item It defines a governed memory lifecycle in which origin determines initial trust and status, corrections supersede rather than overwrite, deletion purges payloads and emits a receipt, and recall considers only active records.
  \item It combines candidate-capped lexical retrieval with an optional bounded graph traversal. The graph is a provenance-preserving, rebuildable index rather than a second source of truth.
  \item It separates extraction confidence, source trust, retrieval score, entity-resolution confidence, and action outcome. None of these numbers silently becomes authority.
  \item It binds memory, graph, retrieval, and effect transitions into one per-subject hash chain with externally anchorable checkpoints and crash-aware action states. The evaluated snapshot adds a versioned digest envelope for selected retrieval and record fields; it does not claim to authenticate every database column.
  \item It describes the implemented platform and filesystem security boundaries precisely, including important non-guarantees.
\end{enumerate}

The claim is architectural, not a claim of priority or universal superiority. The individual mechanisms---full-text search, graph traversal, provenance, hash chains, approval tokens, and compensating actions---are established ideas. The systems contribution is their composition around an explicit distinction between \emph{information} and \emph{authority}, together with a local implementation whose transitions can be inspected and replayed.

\subsection{How to read this paper}

Non-specialist readers can follow the plain-language boxes, \cref{fig:architecture}, the lifecycle table in \cref{tab:lifecycle}, the forensics results in \cref{tab:forensics}, and the platform matrix in \cref{tab:platform-encryption}. Technical readers will find the scoring equations, graph bounds, transaction model, cryptographic construction, forensics protocol, action state machine, implementation map, and reproducibility details in the main text and appendix.

\subsection{Terminology: is this RAG?}

Lewis et al. introduced retrieval-augmented generation (RAG) as a generator combining parametric memory with a dense vector index accessed by a neural retriever~\cite{lewis2020rag}. The term is now also used more broadly for almost any process that inserts retrieved material into a model prompt. Under that broad usage, \system{} recall can be an augmentation source. Under the original architectural meaning, \system{} is not a conventional dense RAG pipeline: it has no embedding model or vector index, and its primary responsibility is governed persistent state and effects rather than passage selection for generation.

We therefore use the precise phrase \emph{governed memory and effect-control plane}. This does not deny that an application can use its returned records in a RAG-like prompting step. It states that retrieval augmentation is an integration point, not an adequate description of the system.
